\section{光学設計と物理的背景}

\subsection{結像系の選定：ピンホールカメラの採用理由}
\subsubsection{設計思想 (Design Philosophy)}
本装置では、軟X線（SXR）の結像方式として「ピンホールカメラ方式」を採用した。
高度な集光光学系（ミラー等）ではなく、あえて単純なピンホールを選んだ理由は、本研究特有の以下の要求仕様に基づく。

\begin{enumerate}
    \item 広視野 (Wide FOV) の確保:
    プラズマの断面全体（数10cm四方）を一度に捉え、磁気プローブ等の障害物を避けて観測する必要がある\cite{Xiang2021}。視野の狭いミラー光学系では、このグローバルな観測が困難である。
    
    \item エネルギー弁別のための波長無依存性:
    本計測の核心は「多色計測（Multi-color）」によるエネルギー解析にある。反射率がX線エネルギーに強く依存するミラー光学系とは異なり、ピンホールは幾何光学的に結像するため、焦点距離や結像特性が光子エネルギーに依存しない。これにより、フィルター交換のみで純粋なエネルギー比較が可能となる。
\end{enumerate}

\subsubsection{ピンホール径 $\phi 1\,\mathrm{mm}$ の決定ロジック}
ピンホール径 $d$ は、要求される「空間分解能」と「信号強度」のトレードオフによって決定された。

\begin{itemize}
    \item 要求分解能: トモグラフィ解析におけるグリッドサイズやプラズマの主要構造を考慮すると、本装置の目標空間分解能は 約$4\,\mathrm{mm}$ である\cite{Xiang2021}。
    \item 信号強度の制約: 観測対象である $100\,\mathrm{eV}$ 以上の軟X線は、プラズマからの放射総量のごく一部であるため、信号は極めて微弱である。
    \item 結論: 幾何光学的なボケが目標分解能（$4\,\mathrm{mm}$）に対して無視できる範囲（$< 1\,\mathrm{mm}$）で、可能な限り開口を大きくして光量を稼ぐため、径を $1\,\mathrm{mm}$ に設定した\cite{Xiang2021}。回折限界の影響はこの領域では無視できる。
\end{itemize}

\subsection{検出器の選定：MCPの動作原理と必然性}
\subsubsection{観測対象：非熱的電子 (Non-thermal Electrons)}
本装置が捉えようとしているのは、マクスウェル分布（熱平衡状態）に従う電子の「裾野（テール）」ではない。
TS-6プラズマのバルク電子温度は $T_e \approx 10\,\mathrm{eV}$ 程度である\cite{Takeda2023}。
したがって、マクスウェル分布を仮定した場合、$100\,\mathrm{eV}$（$10 T_e$）以上のエネルギーを持つ電子は確率的にほぼ存在しない。
すなわち、フィルター越しに $100\,\mathrm{eV}$ 以上の軟X線が観測された場合、それは磁気リコネクション電場等によって加速された「非熱的（Non-thermal）電子」からの制動放射であると断定できる\cite{Takeda2023}。

\subsubsection{なぜフォトダイオードやCCDではダメなのか（重要）}
「光量さえ十分にあれば、半導体センサーでも良いのではないか？」という疑問に対し、本装置においては明確にNoである。
単なる感度不足以外に、以下の物理的・技術的な障壁が存在するため、MCP以外の選択肢はあり得ない。

\begin{enumerate}
    \item 表面不感層（Dead Layer）の問題:
    一般的なフォトダイオードやCCDの表面には保護膜（SiO$_2$等）や電極層が存在する。$100\,\mathrm{eV}$ 程度の軟X線は透過力が極めて弱く、物質表面の数nm〜数十nmで吸収されてしまう。そのため、センサーの感応領域（空乏層）に届く前に表面で信号が消失し、「光は来ているのに感度ゼロ」という状態になる。これを回避するための裏面照射型や窓なしセンサーは極めて高価かつ取り扱いが困難である。
    
    \item 読み出し速度の限界:
    本計測では $5\,\mu\mathrm{s}$ 間隔（$200\,\mathrm{kHz}$）での連続撮影が必要である。$100 \times 100$ 画素程度の2次元画像をこの速度で読み出せるX線直接検知型センサーは、現時点で入手困難な特殊品に限られる。
    
    \item 真空配線の複雑化:
    多チャンネルの半導体アレイを真空中に設置する場合、真空フィードスルーを通して数千〜数万本の信号線を取り出すか、真空内で高度な集積回路を動作させる必要がある。これはノイズ対策や冷却の観点から現実的ではない。
\end{enumerate}

\subsubsection{MCPによる解決策}
以上の問題を、MCP (Micro-Channel Plate) は以下のように解決している。
\begin{itemize}
    \item 表面検知: チャンネル内壁に入射した軟X線が直接電子を叩き出すため、デッドレイヤーの問題がない。
    \item 超高速応答: 応答速度は電子の走行時間と蛍光体の残光時間（P46なら $0.2 \mu\mathrm{s}$）のみで決まるため、MHzオーダーの現象も追随可能。
    \item 光伝送: 信号を「光」に変換してファイバーで真空外に出すため、配線は高圧ケーブル数本で済み、ノイズにも極めて強い。
\end{itemize}

\subsection{像転送光学系：光量不足との戦い}
MCP蛍光面に出力された像を、真空容器外のカメラまで届ける必要がある。
初期設計では光量不足によりS/N比が悪化する問題が発生したため、現在は以下の「高NA（開口数）・拡大」光学系へ変更されている\cite{Takeda2023}。

\subsubsection{1. 真空側レンズ (S-mount Lens)}
\begin{itemize}
    \item 選定: 焦点距離 $f=2.5\,\mathrm{mm}$、F値 F2.0 の明るいレンズを採用。
    \item 役割: MCP蛍光面（$\phi 17\,\mathrm{mm}$）の像を、後段のファイバーバンドル端面に効率よく集光・縮小投影する。F値の暗いレンズは光量ロスの主因となるため不可。
\end{itemize}

\subsubsection{2. 光ファイバーバンドル (Light Guide)}
\begin{itemize}
    \item 仕様: 多成分ガラスファイバーバンドル (6000 fibers/bundle, NA=0.45)。
    \item 物理: $NA=0.45$（受光角 約$27^\circ$）という高い開口数を持つファイバーを選定することで、前段の明るいレンズからの広角な入射光を漏らさず転送する。
\end{itemize}

\subsubsection{3. カメラ側レンズ (LAOWA Ultra-Macro)}
ファイバー出口の像は非常に小さい。これを一般的なレンズ（倍率1倍以下）で撮影すると、カメラセンサの面積を有効活用できず、解像度と光量が無駄になる。
\begin{itemize}
    \item 解決策: LAOWA 60mm F2.8 2X Ultra-Macro を採用。
    \item 効果: 最大2倍の拡大撮影が可能であるため、ファイバー端面の小さな像をCMOSセンサ（$13\,\mathrm{mm} \times 13\,\mathrm{mm}$）一杯に引き伸ばして投影できる。これにより、実効的な画素数を稼ぎ、微弱な信号を余すことなくデジタル化することに成功した\cite{Takeda2023}。
\end{itemize}